% !PRINTED 677
\weberpage{677}
Dies beweist nun der Satz:
% [source note] The source prints der Satz; den Satz is apparently intended, since the preceding argument establishes the theorem which follows. The printed wording reverses that logical direction if read literally.

4. \gesperrt{Ist $g_1, g_2, \ldots, g_\nu$ ein beliebiges System reeller Zahlen, die nicht alle einander gleich sind, so l\"asst sich in jedem algebraischen K\"orper $\Omega$ eine Einheit $\varepsilon$ so bestimmen, dass die Summe}
\[
g_1 l_1 + g_2 l_2 + \cdots + g_\nu l_\nu
\]
\gesperrt{von Null verschieden ist, wenn $l_s$ das System der conjugirten Logarithmen von $\varepsilon$ ist.}

Dieser wichtige Satz, der das Fundament f\"ur das Folgende ist, r\"uhrt, wenn auch in etwas ver\"anderter Fassung, von \gesperrt{Dirichlet} her. Diese Formulirung ist von \gesperrt{Minkowski} gegeben.

\begin{center}
\S.\,186.\\
\gesperrt{Systeme unabh\"angiger Einheiten und Exponentensysteme der Einheiten.}
\end{center}

Wenn wir in dem zuletzt bewiesenen Satze $g_1 = 1, g_2 = 0, \ldots, g_\nu = 0$ setzen, was, wenn wir von jetzt an $\nu > 1$ voraussetzen, gestattet ist, so folgt, dass es in $\Omega$ eine Einheit $\varepsilon$ giebt, f\"ur die einer der conjugirten Logarithmen, $l_1$, von Null verschieden ist. Da jede Potenz von $\varepsilon$ dieselbe Eigenschaft hat, so giebt es auch unendlich viele solche Einheiten.

Dieser Satz gestattet nun eine sehr wichtige Verallgemeinerung:

5. \gesperrt{Ist $s \leqq \nu - 1$, so l\"asst sich im K\"orper $\Omega$ ein System von Einheiten mit den zugeh\"origen conjugirten Logarithmen}
\weq{1}{
\begin{array}{l}
\varepsilon_1 : l_{1,1}, l_{1,2}, \ldots, l_{1,\nu} \\
\varepsilon_2 : l_{2,1}, l_{2,2}, \ldots, l_{2,\nu} \\
\cdots\cdots\cdots\cdots\cdots\cdots \\
\varepsilon_s : l_{s,1}, l_{s,2}, \ldots, l_{s,\nu}
\end{array}
}
\gesperrt{derart bestimmen, dass eine beliebige der $s$-reihigen Determinanten der Matrix der $l_{h,k}$, etwa}
\[
L_s = \Sigma \pm l_{1,1} l_{2,2} \ldots l_{s,s}
\]
\gesperrt{von Null verschieden ist.}

F\"ur $s = 1$ ist der ausgesprochene Satz nach der am Eingange gemachten Bemerkung richtig. Wir nehmen also an, er
%CONT>>
